\documentclass[11pt,a4paper]{article}

% ---------------------------------------------------------
% PREAMBLE
% ---------------------------------------------------------
\usepackage[utf8]{inputenc}
\usepackage[T1]{fontenc}
\usepackage{lmodern}
\usepackage{amsmath,amssymb,amsfonts,bm}
\usepackage{geometry}
\usepackage{graphicx}
\usepackage{hyperref}
\usepackage{cite}
\usepackage{authblk}
\usepackage{physics}
\usepackage{mathtools}
\usepackage{enumitem}

\geometry{margin=2.4cm}

\title{\textbf{Companion XLIV (v49.1)}\\[0.2cm]
\textbf{Gravitational Slip in the Relaxation-Driven Cyclic Cosmology}\\[0.2cm]
\large Foundations of Relaxation-Driven Cyclic Cosmology (RDCC)}

\author[1]{Michael Lehmann}
\affil[1]{\small Independent Researcher, Relaxation-Driven Cyclic Cosmology (RDCC)\\
\texttt{mi.lehmann@gmx.de}}

\date{September 2026}

\begin{document}
\maketitle

\begin{abstract}
The gravitational slip provides one of the most sensitive tests of modifications
to the geometric structure of spacetime. In previous versions of RDCC, the slip
was described through phenomenological scale-dependent modifications of the
scalar metric potentials. In the present Phase-1 integrated version, the
gravitational slip is derived from an effective anisotropic stress generated by
the combined action of the mirror sector and the dimension-8 inter-sector
operator. The resulting framework establishes a direct causal chain between the
global relaxation dynamics, the freeze-out geometry, the relaxation scale
\(k_{\mathrm{rel}}\), and the observable relation between the metric potentials
\(\Phi\) and \(\Psi\). This Companion provides the geometric foundation for weak
lensing, redshift-space distortions, integrated Sachs--Wolfe measurements, and
future survey constraints on RDCC.
\end{abstract}

\tableofcontents
\newpage

% ---------------------------------------------------------
\section{Introduction}
% ---------------------------------------------------------

In the Newtonian gauge, scalar perturbations of the metric are described by two
potentials, \(\Phi\) and \(\Psi\), which govern the motion of matter and the
propagation of light. In the standard cosmological model these potentials are
equal in the absence of anisotropic stress, and their evolution is governed by
the linear growth of structure.

In earlier versions of RDCC the slip was introduced by scale-dependent
rescalings of the two potentials. That description remains useful as a fitting
template, but it is no longer the primary definition. In v49.1 the slip is
generated geometrically: the mirror sector and the dimension-8 operator source
an effective anisotropic stress, and Companion XII supplies the freeze-out scale
\(k_{\mathrm{rel}}\) at which this stress becomes important. The slip is then a
tracer of that same geometric transition.

% ---------------------------------------------------------
\section{Metric Potentials in RDCC}
% ---------------------------------------------------------

Scalar perturbations are described in the Newtonian gauge by
\begin{equation}
  ds^{2}
  =
  -(1+2\Psi)\,dt^{2}
  +
  a^{2}(t)\,(1-2\Phi)\,d\mathbf{x}^{2}.
\end{equation}
The potentials \(\Phi\) and \(\Psi\) govern the motion of matter and the
propagation of light. In standard cosmology, the two potentials remain equal in
the absence of anisotropic stress,
\begin{equation}
  \Phi = \Psi.
\end{equation}
Within RDCC, departures from this equality arise naturally from the geometric
structure of the mirror sector and the dimension-8 relaxation operator.

% ---------------------------------------------------------
\section{Effective Anisotropic Stress}
% ---------------------------------------------------------

The leading interaction between the visible and mirror sectors is described by
\begin{equation}
  \mathcal{O}_{\mathrm{int}}
  =
  \frac{1}{M^{4}}
  T_{\mu\nu}^{(+)} T^{(-)\mu\nu}.
\end{equation}
The variation of this operator generates an effective geometric contribution to
the stress-energy tensor. At the perturbative level, this contribution acts as
an effective anisotropic stress.

We define the total effective stress as
\begin{equation}
  \pi_{\mathrm{eff}}
  =
  \pi_{\mathrm{mirror}}
  +
  \pi_{\mathrm{rel}},
\end{equation}
where \(\pi_{\mathrm{mirror}}\) arises from the diffuse mirror-sector geometry
and \(\pi_{\mathrm{rel}}\) encodes the response of the relaxation sector.

The quantity \(\pi_{\mathrm{mirror}}\) does not represent a conventional matter
anisotropic stress. Instead, it arises effectively from the diffuse geometric
response of the mirror sector.

% ---------------------------------------------------------
\section{Origin of the Gravitational Slip}
% ---------------------------------------------------------

The scalar Einstein equations imply
\begin{equation}
  k^{2}
  \left(
    \Phi - \Psi
  \right)
  =
  8\pi G a^{2}\,
  \pi_{\mathrm{eff}}.
  \label{eq:sliporigin}
\end{equation}
Equation~(\ref{eq:sliporigin}) shows that the gravitational slip is not
introduced phenomenologically. Instead, it emerges from the effective
anisotropic stress generated by the mirror sector and the relaxation geometry.
The gravitational slip therefore becomes a direct geometric consequence of the
global CPT-symmetric structure of RDCC.

% ---------------------------------------------------------
\section{Connection to the Relaxation Scale}
% ---------------------------------------------------------

Companion XII established the relaxation scale
\begin{equation}
  k_{\mathrm{rel}}
  =
  a_{\mathrm{fr}} H(a_{\mathrm{fr}}).
\end{equation}
This scale separates perturbations that evolve almost identically to their
standard cosmological counterparts from perturbations affected by relaxation
dynamics.

The effective anisotropic stress becomes increasingly important for
\begin{equation}
  k \gtrsim k_{\mathrm{rel}}.
\end{equation}
Consequently, the largest deviations from
\begin{equation}
  \Phi = \Psi
\end{equation}
are expected near and above the relaxation scale.

% ---------------------------------------------------------
\section{Definition of the Gravitational Slip}
% ---------------------------------------------------------

The slip parameter is defined as
\begin{equation}
  \eta(k,a)
  =
  \frac{\Phi(k,a)}{\Psi(k,a)}.
\end{equation}
In the absence of effective anisotropic stress,
\begin{equation}
  \eta = 1.
\end{equation}
Within RDCC,
\begin{equation}
  \eta \neq 1
\end{equation}
arises dynamically through Eq.~(\ref{eq:sliporigin}). The scale dependence of
\(\eta\) is therefore inherited from the relaxation scale and the perturbative
structure of the mirror sector.

The older phenomenological template
\begin{equation}
  \eta_{\chi}(k,a)
  =
  \frac{
    1 - \chi_{\Phi}
    \left(k/k_{\mathrm{rel}}\right)^{\alpha}
  }{
    1 - \chi_{\Psi}
    \left(k/k_{\mathrm{rel}}\right)^{\alpha}
  }
\end{equation}
may still be used in survey fits. It is not the primary definition and must
remain consistent with Eq.~(\ref{eq:sliporigin}).

% ---------------------------------------------------------
\section{Observational Signatures}
% ---------------------------------------------------------

The gravitational slip affects several key observables. Weak lensing is
sensitive to the combination \(\Phi+\Psi\), while galaxy motions are sensitive
to \(\Psi\) alone. The slip therefore modifies the relation between lensing and
clustering.

The integrated Sachs--Wolfe effect is sensitive to the time derivative of the
metric potentials, and the slip introduces a scale-dependent modification to
the ISW signal.

Redshift-space distortions depend on the velocity field, which is governed by
\(\Psi\). The slip therefore introduces a scale-dependent modification to the
RSD signal.

These signatures provide a direct observational window into the relaxation
scale and the geometric structure of the RDCC Universe.

% ---------------------------------------------------------
\section{Conclusion}
% ---------------------------------------------------------

The gravitational slip of RDCC is reinterpreted as a geometric consequence of
the combined mirror-sector and relaxation dynamics. The fundamental causal
chain is
\begin{equation}
  \alpha
  \;\rightarrow\;
  a_{\mathrm{fr}}
  \;\rightarrow\;
  k_{\mathrm{rel}}
  \;\rightarrow\;
  \pi_{\mathrm{eff}}
  \;\rightarrow\;
  (\Phi-\Psi).
\end{equation}
The slip is therefore no longer a phenomenological assumption but an emergent
consequence of the global CPT-symmetric geometry.

The perturbative source of the slip is ultimately encoded in the mirror-sector
fluctuations \(\delta_{(-)}\) developed in Companion VII.

This framework provides the theoretical basis for future analyses of weak
lensing, cosmic shear, redshift-space distortions, and integrated Sachs--Wolfe
observables within RDCC.

% ========================
% References
% ========================
\begin{thebibliography}{10}

\bibitem{Flagship} Lehmann, M. (2026). Relaxation-Driven Cyclic Cosmology
(RDCC): A Minimal CPT-Symmetric Framework with a Single Parameter. Flagship
Paper. Zenodo: \url{https://zenodo.org/records/19744958} (v43.0 or later).

\bibitem{Zenodo} The complete RDCC companion series is available at the same
Zenodo record above.

\bibitem{CompVII} Companion VII (v49.1): Boltzmann Hierarchy in a CPT-Symmetric
Bipartite Universe.

\bibitem{CompXII} Companion XII (v49.1): Relaxation Scale and Freeze-Out
Geometry.

\bibitem{CompXI} Companion XI (v49.1): Background Cosmology and Modified
Friedmann Equations.

\bibitem{Planck2020} Planck Collaboration (2020), Astron.\ Astrophys.\ 641, A6.

\bibitem{DESI2024} DESI Collaboration (2024/2025), BAO and \(f\sigma_8\) analyses.

\bibitem{Kaiser1987} Kaiser, N. (1987), Mon.\ Not.\ R.\ Astron.\ Soc.\ 227, 1.

\bibitem{Bartelmann2001} Bartelmann, M.\ \& Schneider, P. (2001), Phys.\ Rep.\
340, 291.

\bibitem{HuOkamoto2002} Hu, W.\ \& Okamoto, T. (2002), Astrophys.\ J.\ 574, 566.

\bibitem{LewisChallinor2006} Lewis, A.\ \& Challinor, A. (2006), Phys.\ Rep.\
429, 1.

\end{thebibliography}

\end{document}